\documentclass[11pt,a4paper]{article}

\usepackage[utf8]{inputenc}
\usepackage[T1]{fontenc}
\usepackage[english]{babel}
\usepackage{amsmath,amssymb,amsthm,mathtools,bm}
\usepackage[a4paper,margin=2.5cm]{geometry}
\usepackage{hyperref}
\usepackage{booktabs}
\usepackage{array}
\usepackage{xcolor}
\usepackage{enumitem}

\hypersetup{
  colorlinks=true,
  linkcolor=blue!50!black,
  citecolor=blue!50!black,
  urlcolor=blue!50!black
}

\theoremstyle{definition}
\newtheorem{definition}{Definition}[section]
\newtheorem{algorithm}[definition]{Algorithm}
\newtheorem{example}[definition]{Example}

\theoremstyle{plain}
\newtheorem{theorem}[definition]{Theorem}
\newtheorem{lemma}[definition]{Lemma}
\newtheorem{proposition}[definition]{Proposition}
\newtheorem{corollary}[definition]{Corollary}

\theoremstyle{remark}
\newtheorem{remark}[definition]{Remark}

\newcommand{\C}{\mathbb C}
\newcommand{\R}{\mathbb R}
\newcommand{\N}{\mathbb N}
\newcommand{\eps}{\varepsilon}
\newcommand{\dd}{\,\mathrm d}
\newcommand{\Res}{\operatorname{Res}}
\newcommand{\diag}{\operatorname{diag}}
\newcommand{\Id}{\operatorname{Id}}
\newcommand{\norm}[1]{\left\lVert #1\right\rVert}
\newcommand{\abs}[1]{\left\lvert #1\right\rvert}

\title{\textbf{Pandrosion Geometry in Several Complex Variables}\\[0.35em]
\large Integral slopes, homothetic transport, guarded continuation,\\
and the resultant--projective completion architecture of Flow 092}
\author{Ivan Besevic}
\date{April 2026}

\begin{document}

\maketitle

\begin{abstract}
We formulate the multidimensional Pandrosion construction used in the
implementation \texttt{flow/092\_pandrosion\_resultant\_pairing.py}.  The
setting is a square polynomial system
\[
  F=(F_1,\ldots,F_n):\C^n\longrightarrow\C^n
\]
with isolated regular zeros.  The basic object is not a derivative at a point,
but a two-point slope matrix \(Q(a,z)\) satisfying the exact telescoping
identity
\[
  F(z)-F(a)=Q(a,z)(z-a).
\]
The canonical geometric representative is the segment-integral matrix
\[
  Q^{\mathrm{int}}_{ij}(a,z)
  =
  \int_0^1
  \partial_jF_i\bigl(a+t(z-a)\bigr)\,\dd t,
\]
which extends Pandrosion's one-dimensional divided quotient to arbitrary
dimension.  We then describe the system-adaptive exactified slope, convex
blends of exact geometries, diagonal homothetic transport, Aitken--Pandrosion
\(T_2\) acceleration, and the branch-safe continuation guard.

The resulting algorithm remains a homotopy method on \(\C^n\), but the
predictor--corrector step is derivative-free in the Pandrosion sense: every
accepted candidate is corrected by a non-local slope matrix rather than by a
Newton--Lairez corrector.  Flow 092 adds three completion layers: coordinate
resultants for bivariate systems, deterministic multivariate gamma sweeps, and
projective/Riemann-sphere charts.  Resultants and projective charts are used
only as candidate generators; acceptance is always by Pandrosion polish on the
original system.  We record the mathematical identities behind the construction,
state local fixed-point and transport propositions, and give an empirical
summary of the current default configuration.  No complexity-theoretic claim
against the Beltran--Pardo or Lairez algorithms is made.
\end{abstract}

\paragraph{Implementation note.}
The paper describes the construction implemented by the chain
\[
092\to091\to090\to089\to087\to081\to076\to070\to069\to068.
\]
The numbering is repository-local.  The mathematical object is the
multidimensional Pandrosion slope geometry; the code numbers are included only
to make the implementation auditable.

\tableofcontents

%==============================================================================
\section{Introduction}
%==============================================================================

Pandrosion's original construction for \(p\)-th roots can be written as a
fixed-point iteration generated by a geometric sum.  In one complex variable,
the essential identity is the divided quotient
\[
  P(z)-P(a)=Q_P(a,z)(z-a),
  \qquad
  Q_P(a,z)=\frac{P(z)-P(a)}{z-a}.
\]
If \(\zeta\) is a zero of \(P\), then
\[
  \zeta
  =
  a-Q_P(a,\zeta)^{-1}P(a),
\]
which suggests a root-refinement map
\[
  z\longmapsto a-Q_P(a,z)^{-1}P(a).
\]
This map is not Newton's method unless the anchor \(a\) is collapsed onto the
current iterate.  The non-locality of \(Q_P(a,z)\) is the geometric content of
Pandrosion: the step is built from a chord, not a tangent.

For a system \(F:\C^n\to\C^n\), the scalar quotient is replaced by a slope
matrix \(Q_F(a,z)\).  The present paper develops the form used in the 092
solver: a portfolio of exact two-point geometries, transported by a diagonal
homothety and embedded in a total-degree homotopy.  The root-finding problem is
the standard one:
\[
  F(z)=0,\qquad z\in\C^n.
\]
All variables, coefficients, homotopy parameters, and candidate roots live over
\(\C\).

\paragraph{What is new in 092 relative to the earlier flows.}
The earlier flows establish the Pandrosion corrector and the homothetic
geometry.  Flow 092 adds a practical completion architecture:
\begin{itemize}[leftmargin=2em]
  \item branch-safe path tracking with residual and predictor-compatibility
  guards;
  \item bivariate coordinate-resultant recovery when \(n=2\);
  \item deterministic multivariate completion sweeps when \(n\ge3\);
  \item projective/Riemann-sphere recovery charts as a last-resort tool;
  \item tight bivariate clustering, because valid roots can be closer than the
  generic \(10^{-6}\) clustering radius.
\end{itemize}
The empirical lesson of the current implementation is that these layers are
most effective in this order: homothetic Pandrosion tracking first, then
resultant/projective recovery only for missing sheets.

%==============================================================================
\section{Polynomial systems over \texorpdfstring{\(\C\)}{C}}
%==============================================================================

Let
\[
  F=(F_1,\ldots,F_n),\qquad F_i\in\C[z_1,\ldots,z_n],
\]
be a square system.  We write \(a,z\in\C^n\), \(\Delta=z-a\), and
\[
  J_F(x)=\left(\partial_jF_i(x)\right)_{1\le i,j\le n}.
\]
The B\'ezout number of the system is
\[
  D=\prod_{i=1}^n d_i,\qquad d_i=\deg F_i.
\]
For the families used in the experiments, \(D\) is the target number of isolated
solutions counted with multiplicity.

\begin{definition}[Pandrosion slope matrix]
A matrix-valued map \(Q(a,z)\in\C^{n\times n}\) is a Pandrosion slope for
\(F\) between \(a\) and \(z\) if
\begin{equation}
  \label{eq:telescope}
  F(z)-F(a)=Q(a,z)(z-a).
\end{equation}
We call \eqref{eq:telescope} the telescoping identity.
\end{definition}

The identity is the load-bearing invariant of the construction.  Different
geometries correspond to different choices of \(Q\), but all admissible choices
must satisfy \eqref{eq:telescope}.

\begin{definition}[Pandrosion correction map]
Let \(Q(a,z)\) be an invertible slope matrix.  The associated anchored
Pandrosion map is
\begin{equation}
  \label{eq:pandrosion-map}
  \mathcal P_Q(a,z)
  =
  a-Q(a,z)^{-1}F(a).
\end{equation}
\end{definition}

\begin{proposition}[Fixed-zero property]
\label{prop:fixed-zero}
If \(\zeta\) is a zero of \(F\) and \(Q(a,\zeta)\) is invertible, then
\[
  \mathcal P_Q(a,\zeta)=\zeta.
\]
\end{proposition}

\begin{proof}
Since \(F(\zeta)=0\), the telescoping identity gives
\[
  -F(a)=Q(a,\zeta)(\zeta-a).
\]
Multiplying by \(Q(a,\zeta)^{-1}\) yields
\[
  -Q(a,\zeta)^{-1}F(a)=\zeta-a,
\]
which is equivalent to \(\mathcal P_Q(a,\zeta)=\zeta\).
\end{proof}

\begin{remark}
Proposition~\ref{prop:fixed-zero} is the multidimensional analogue of the
classical one-dimensional Pandrosion identity.  It explains why a two-point
slope can act as a corrector: at an actual zero, the corrector fixes the zero
exactly.
\end{remark}

%==============================================================================
\section{The integral geometry}
%==============================================================================

The cleanest slope matrix is obtained by integrating the Jacobian along the
straight segment from \(a\) to \(z\).

\begin{definition}[Segment-integral Pandrosion slope]
\label{def:integral-slope}
For \(a,z\in\C^n\), define
\begin{equation}
  \label{eq:q-int}
  Q^{\mathrm{int}}_{ij}(a,z)
  =
  \int_0^1
  \partial_jF_i\bigl(a+t(z-a)\bigr)\,\dd t.
\end{equation}
Equivalently,
\[
  Q^{\mathrm{int}}(a,z)
  =
  \int_0^1 J_F(a+t(z-a))\,\dd t.
\]
\end{definition}

\begin{theorem}[Exact telescoping for \(Q^{\mathrm{int}}\)]
\label{thm:int-telescope}
For every polynomial map \(F:\C^n\to\C^n\),
\[
  F(z)-F(a)=Q^{\mathrm{int}}(a,z)(z-a).
\]
\end{theorem}

\begin{proof}
Let \(\gamma(t)=a+t(z-a)\).  For each component \(F_i\),
\[
  \frac{\dd}{\dd t}F_i(\gamma(t))
  =
  \sum_{j=1}^n
  \partial_jF_i(\gamma(t))(z_j-a_j).
\]
Integrating from \(0\) to \(1\) gives
\[
  F_i(z)-F_i(a)
  =
  \sum_{j=1}^n
  \left(\int_0^1\partial_jF_i(a+t(z-a))\,\dd t\right)(z_j-a_j),
\]
which is the \(i\)-th row of the matrix identity.
\end{proof}

\paragraph{Numerical realization.}
For polynomial systems, the integrand in \eqref{eq:q-int} is polynomial in
\(t\).  A Gauss--Legendre rule of sufficient order is exact.  The implementation
uses a capped exact/approximate Gauss--Legendre rule:
\[
  Q^{\mathrm{GL}}(a,z)
  =
  \sum_{\ell=1}^m w_\ell J_F(a+t_\ell(z-a)).
\]
When the cap is high enough relative to \(\max_i d_i\), this equals
\(Q^{\mathrm{int}}\) up to roundoff.

\begin{remark}[Classical interpretation]
The matrix \(Q^{\mathrm{int}}\) is a Schmidt-type divided-difference matrix for
systems.  The Pandrosion contribution here is not the existence of divided
differences; it is the use of these matrices as geometric two-point correctors,
combined with homothetic scaling and guarded continuation.
\end{remark}

%==============================================================================
\section{System-adaptive exactification}
%==============================================================================

The integral geometry is robust but can be expensive.  Flow 070 introduced a
cheaper system-adaptive geometry.  It begins from an edge slope frame \(B(a,z)\)
and then performs a rank-one exactification to restore the telescoping identity.

\subsection{Edge slope frame}

Let \(e^{(j)}\) be the point obtained from \(a\) by replacing only the
\(j\)-th coordinate with \(z_j\).  Define
\[
  B_{ij}(a,z)
  =
  \begin{cases}
  \dfrac{F_i(e^{(j)})-F_i(a)}{z_j-a_j},
  & z_j\ne a_j,\\[1.2em]
  0,& z_j=a_j.
  \end{cases}
\]
This matrix is cheap and locally informative, but it does not generally satisfy
\eqref{eq:telescope}.

\subsection{Rank-one exactification}

Let
\[
  r(a,z)=F(z)-F(a)-B(a,z)(z-a)
\]
be the defect of the edge frame.  Choose a covector \(w(a,z)\in(\C^n)^*\) such
that \(w(a,z)(z-a)\ne0\).  Define
\begin{equation}
  \label{eq:exactify}
  Q^{\mathrm{sys}}(a,z)
  =
  B(a,z)
  +
  \frac{r(a,z)\,w(a,z)}{w(a,z)(z-a)}.
\end{equation}

\begin{proposition}[Exactification]
\label{prop:exactify}
The matrix \(Q^{\mathrm{sys}}\) defined by \eqref{eq:exactify} satisfies
\[
  F(z)-F(a)=Q^{\mathrm{sys}}(a,z)(z-a).
\]
\end{proposition}

\begin{proof}
Multiplying \eqref{eq:exactify} by \(z-a\) gives
\[
  Q^{\mathrm{sys}}(a,z)(z-a)
  =
  B(a,z)(z-a)
  +
  r(a,z)
  =
  F(z)-F(a).
\]
\end{proof}

\paragraph{The adaptive covector.}
In the implementation, \(w(a,z)\) is chosen from the actual system:
support activity, coefficient magnitudes, local monomial energy on the box
joining \(a\) and \(z\), midpoint Jacobian column energy, and residual
alignment.  This does not change the proof of Proposition~\ref{prop:exactify};
it changes only the conditioning of the resulting matrix.

\subsection{Blended geometries}

If \(Q_1\) and \(Q_2\) both satisfy \eqref{eq:telescope}, then so does
\[
  Q_\theta=\theta Q_1+(1-\theta)Q_2
  \qquad(0\le \theta\le1).
\]
Flow 092 inherits the mode portfolio
\[
  \mathrm{system},\qquad \mathrm{integral\_gl},\qquad \mathrm{blend}.
\]
The blend mode uses the exactified system geometry and the integral geometry.
It is still an exact Pandrosion geometry because exactness is affine.

%==============================================================================
\section{Local acceleration and damping}
%==============================================================================

For a fixed anchor \(a\), define
\[
  H(z)=\mathcal P_Q(a,z).
\]
The implemented corrector first tries the plain Pandrosion step \(H(z)\).  If
the residual decreases sufficiently, it attempts a coordinatewise Aitken
acceleration.

\begin{definition}[Pandrosion \(T_2\) step]
Let
\[
  s_1=H(z),\qquad s_2=H(s_1).
\]
The \(T_2\) candidate is the vector whose \(k\)-th component is
\begin{equation}
  \label{eq:t2}
  (T_2(z))_k
  =
  z_k-\frac{(s_{1,k}-z_k)^2}{s_{2,k}-2s_{1,k}+z_k},
\end{equation}
with the convention that \(s_{2,k}\) is used when the denominator is too small.
\end{definition}

The step is guarded by two rules:
\begin{itemize}[leftmargin=2em]
  \item if the accelerated point is non-finite or too far from the current
  scale, fall back to \(s_2\);
  \item accept only a damped point
  \[
    z+\tau(c-z),\qquad
    \tau\in\{1,0.72,0.5,0.35,0.25,0.125,0.0625\},
  \]
  that reduces the residual.
\end{itemize}

\begin{remark}[Why \(T_2\), not \(T_3\) or \(T_4\)]
The repository contains an experimental wrapper comparing \(T_2,T_3,T_4\).  In
the tested regimes, \(T_3\) and \(T_4\) sometimes reduce the number of candidate
paths but generally increase wall-clock time, and \(T_4\) can destabilize dense
bivariate cases.  Thus Flow 092 uses \(T_2\) as the default local acceleration.
\end{remark}

%==============================================================================
\section{Homothetic transport}
%==============================================================================

The scalar paper uses scaling to replace a large ratio \(x\) by a reduced ratio
near \(1\).  In several variables, Flow 076 implements the analogue as a
diagonal homothety.

\begin{definition}[Diagonal homothety]
Let \(S=\diag(s_1,\ldots,s_n)\) be an invertible diagonal matrix and let
\(D=\diag(d_1,\ldots,d_n)\) be an optional equation normalization.  Define
\[
  G(y)=D\,F(Sy),
  \qquad z=Sy.
\]
\end{definition}

\begin{proposition}[Transport of Pandrosion slopes]
\label{prop:transport}
Assume \(Q_G(a,y)\) satisfies
\[
  G(y)-G(a)=Q_G(a,y)(y-a).
\]
Then
\[
  Q_F(Sa,Sy)=D^{-1}Q_G(a,y)S^{-1}
\]
satisfies
\[
  F(Sy)-F(Sa)=Q_F(Sa,Sy)(Sy-Sa).
\]
\end{proposition}

\begin{proof}
Since \(G(y)=DF(Sy)\),
\[
  D(F(Sy)-F(Sa))=Q_G(a,y)(y-a).
\]
Multiplying by \(D^{-1}\) and using \(y-a=S^{-1}(Sy-Sa)\) gives the claim.
\end{proof}

The scales \(S\) are generated from coefficient and support balancing, and in
the bivariate case can be blended with a Pandrosion polygon/Jensen estimate.
Equation normalization \(D\) equalizes component magnitudes.  The key point is
that scaling changes the geometry of the path without breaking
\eqref{eq:telescope}.

%==============================================================================
\section{Total-degree homotopy in \texorpdfstring{\(\C^n\)}{C n}}
%==============================================================================

Flow 092 is still a homotopy method.  It uses total-degree starts
\[
  G_i(z)=z_i^{d_i}-\phi_i,
\]
where the phases \(\phi_i\in\C^\times\) are deterministic complex phases.  The
start roots are the Cartesian product of the roots of the \(G_i\).  Their number
is exactly \(D=\prod_i d_i\).

For a vector of nonzero complex phases \(\gamma=(\gamma_1,\ldots,\gamma_n)\),
write
\[
  (\Gamma F)_i=\gamma_iF_i.
\]
The tracked homotopy is
\begin{equation}
  \label{eq:homotopy}
  H_t(z)=(1-t)G(z)+t\,\Gamma F(z),
  \qquad 0\le t\le1.
\end{equation}

\paragraph{Predictor.}
The tracker uses a tangent predictor when stable.  Formally, differentiating the
path identity \(H_t(z(t))=0\) gives
\begin{equation*}
  J_{H_t}(z)\dot z+\partial_t H_t(z)=0.
\end{equation*}
The predictor estimates \(z(t+\Delta t)\) from this relation or from a secant
history.

\paragraph{Corrector.}
At the predicted point, the corrector is not a Newton--Lairez corrector.  It is
the Pandrosion corrector of Sections~2--5 applied to the system \(H_{t+\Delta
t}\).

\paragraph{Branch safety.}
A corrected point is accepted only if:
\[
  \norm{H_{t+\Delta t}(z_{\mathrm{corr}})}
  \quad\text{is small}
\]
and the correction is compatible with the predictor move.  If the correction is
too large relative to the local scale or predictor displacement, the step is
declared ambiguous and the time step is reduced.  This is the purpose of the
branch-safe tracker in Flow 087.

%==============================================================================
\section{Bivariate coordinate resultants}
%==============================================================================

For \(n=2\), Flow 089 and Flow 092 use coordinate resultants as recovery
devices.  Let
\[
  F=(F_1,F_2)\in\C[x,y]^2.
\]
Define
\begin{equation}
  \label{eq:coord-res}
  R_x(x)=\Res_y(F_1(x,y),F_2(x,y)),
  \qquad
  R_y(y)=\Res_x(F_1(x,y),F_2(x,y)).
\end{equation}
If \((\xi,\eta)\) is a common zero of \(F_1,F_2\), then
\[
  R_x(\xi)=0,\qquad R_y(\eta)=0.
\]

\paragraph{Reconstruction.}
The implementation samples \(R_x\) and \(R_y\) on circles and reconstructs
their coefficients by a discrete Fourier formula.  It then computes their roots
with a Durand--Kerner iteration.  These roots are not accepted as solutions of
the original system.  They are used only to generate candidate starts:
\[
  x=\xi\quad\Rightarrow\quad
  F_i(\xi,y)=0
\]
for one or both equations, followed by residual ranking and Pandrosion polish
on the full system.

\paragraph{Pairing layer.}
Flow 092 also contains an opt-in layer that cross-pairs roots of \(R_x\) and
\(R_y\).  In the current best default this pair recovery is disabled, because
it is often more expensive than the recovery it provides.  The more important
092 default is bivariate tight clustering: if the user does not specify a
cluster radius and \(n=2\), valid roots are clustered at a much smaller radius
than the generic multivariate default.

%==============================================================================
\section{Multivariate and projective completion}
%==============================================================================

When \(n\ge3\), coordinate resultants in the literal multivariate sense are
large elimination objects.  Flow 092 therefore uses two different completion
mechanisms.

\subsection{Gamma completion sweeps}

Flow 090 reruns deterministic gamma sweeps over the same total-degree start
set, but only after the base pass has failed to reach the B\'ezout count.  Each
candidate root extracted from these sweeps is accepted only if a Pandrosion
polish on the original target system brings the residual below the acceptance
threshold.

\subsection{Projective/Riemann-sphere charts}

Flow 091 treats each coordinate as a Riemann sphere with two affine charts:
\[
  \text{north: } z_j=y_j,
  \qquad
  \text{south: } z_j=\frac{s_j}{y_j}.
\]
For a chosen mask of south-chart coordinates, denominators are cleared
equation-wise so that the transformed system is again polynomial.  A root found
in chart coordinates is mapped back to the original \(z\)-coordinates and then
polished on the original system.

\begin{remark}[Why projective recovery is late]
The projective chart can improve conditioning for roots that are large or badly
placed in the affine chart.  However, choosing a good chart before seeing any
partial root cloud is difficult.  The current architecture therefore uses
projective/Riemann charts as a recovery tool rather than as the default
starting geometry.
\end{remark}

%==============================================================================
\section{The Flow 092 algorithm}
%==============================================================================

We can now state the algorithm in mathematical form.

\begin{algorithm}[Pandrosion dD solver, Flow 092]
Input: a square polynomial system \(F:\C^n\to\C^n\), degrees \(d_i\), and
B\'ezout target \(D=\prod_i d_i\).

\begin{enumerate}[leftmargin=2em]
  \item Build a diagonal homothety \(z=Sy\), optionally with equation
  normalization \(D_F\).
  \item Build candidate geometry policies:
  system-unit geometry in all dimensions, and for \(n=2\) additional
  polygon/Jensen-informed policies.
  \item Score policies on a small sentinel subset of total-degree paths.
  \item Select one policy and track total-degree homotopy paths in \(\C^n\):
  predictor by tangent/secant extrapolation, corrector by the Pandrosion
  portfolio
  \[
    Q^{\mathrm{sys}},\quad Q^{\mathrm{int}},\quad
    \theta Q^{\mathrm{int}}+(1-\theta)Q^{\mathrm{sys}},
  \]
  with \(T_2\), damping, and branch guard.
  \item If the B\'ezout count is reached, stop.
  \item If \(n=2\), run coordinate-resultant recovery
  \eqref{eq:coord-res}.
  \item If \(n\ge3\), run multivariate gamma completion sweeps.
  \item If roots are still missing, run projective/Riemann-sphere completion.
  \item Optionally, run bivariate pair recovery.
  \item Return the cluster set of candidates whose original residual is below
  the acceptance threshold.
\end{enumerate}
\end{algorithm}

\begin{remark}[Default configuration]
The current 092 defaults are tuned for the observed best configuration:
\[
  \texttt{batch-timeout}=20,\quad
  \texttt{parallel-batches}=4,\quad
  \texttt{batch-size}=16,\quad
  \texttt{micro-batch}=4,
\]
with equation normalization and early stopping at B\'ezout enabled.  Pair
recovery is opt-in.
\end{remark}

%==============================================================================
\section{Correctness statements}
%==============================================================================

The algorithm is heuristic globally, but several of its internal transformations
are exact.

\begin{theorem}[Geometry exactness]
Every slope matrix used by the main Pandrosion corrector in Flow 092 satisfies
the telescoping identity \eqref{eq:telescope}, up to quadrature and floating
roundoff:
\[
  Q^{\mathrm{int}},\qquad Q^{\mathrm{sys}},\qquad
  \theta Q^{\mathrm{int}}+(1-\theta)Q^{\mathrm{sys}}.
\]
\end{theorem}

\begin{proof}
For \(Q^{\mathrm{int}}\), this is Theorem~\ref{thm:int-telescope}.  For
\(Q^{\mathrm{sys}}\), this is Proposition~\ref{prop:exactify}.  The blend is a
convex combination of two exact slopes, hence
\[
  Q_\theta(a,z)(z-a)
  =
  \theta(F(z)-F(a))+(1-\theta)(F(z)-F(a))
  =
  F(z)-F(a).
\]
\end{proof}

\begin{theorem}[Homothety invariance]
Diagonal scaling and equation normalization preserve the class of Pandrosion
slope matrices.
\end{theorem}

\begin{proof}
This is Proposition~\ref{prop:transport}.
\end{proof}

\begin{corollary}[Local fixed-zero property after scaling]
Let \(\zeta\) be a regular zero of \(F\).  In any scaled coordinate chart
\(z=Sy\), the corresponding scaled zero \(S^{-1}\zeta\) is a fixed point of
the scaled Pandrosion map whenever the transported slope matrix is invertible.
\end{corollary}

\begin{proof}
Combine Proposition~\ref{prop:fixed-zero} with Proposition~\ref{prop:transport}.
\end{proof}

\begin{remark}[What is not proved]
This paper does not prove that Flow 092 finds all roots of every polynomial
system, nor does it claim an average polynomial-time bound.  The established
average-case theory for polynomial homotopy over \(\C\) is due to
Beltran--Pardo and Lairez.  The purpose here is to present a derivative-free
Pandrosion geometry and document a strong empirical solver built from it.
\end{remark}

%==============================================================================
\section{Empirical behavior}
%==============================================================================

The following table records representative local runs of the 092 architecture
on the repository's synthetic families.  Timings are machine-dependent and are
included only to indicate scale.

\begin{center}
\begin{tabular}{llrrrrl}
\toprule
Family & case & B\'ezout & roots & paths & seconds & status \\
\midrule
\texttt{sparse\_ks} & \((2,12)\) & 144 & 144 & 201 & \( \approx 4\) & ok \\
\texttt{dense064}   & \((2,10)\) & 100 & 100 & 110 & \( \approx 7\) & ok \\
\texttt{ks}         & \((2,12)\) & 144 & 144 & 232 & \( \approx 35\) & ok \\
\texttt{ks}         & \((3,4)\)  & 64  & 64  & 192 & \( \approx 8\) & ok \\
\texttt{ks}         & \((3,5)\)  & 125 & 125 & 250 & \( \approx 20\) & ok \\
\bottomrule
\end{tabular}
\end{center}

\paragraph{Negative experiments.}
Two natural alternatives were tested and rejected as defaults.
\begin{itemize}[leftmargin=2em]
  \item Higher \(T_n\) acceleration (\(T_3,T_4\)) did not improve wall-clock
  time reliably over \(T_2\).
  \item Resultant-first scheduling was inferior to tracking-first scheduling
  except on very small bivariate cases.  Resultants are best used as recovery,
  not as the main engine.
\end{itemize}

%==============================================================================
\section{Conclusion}
%==============================================================================

The multidimensional Pandrosion geometry is governed by one simple invariant:
\[
  F(z)-F(a)=Q(a,z)(z-a).
\]
The segment-integral slope \(Q^{\mathrm{int}}\) proves that this invariant is
canonical in any dimension; the system-adaptive exactification shows how to
construct cheaper geometries without losing exactness; homothetic transport
preserves the identity under scaling; and the branch-safe homotopy tracker uses
these slopes as practical correctors over \(\C^n\).

Flow 092 is therefore not merely a collection of numerical tricks.  It is a
layered implementation of a coherent geometric principle:
\begin{align*}
  &\text{exact two-point geometry}
  +\text{homothetic conditioning} \\
  &\hspace{2.5cm}
  +\text{guarded continuation}
  +\text{Pandrosion-certified recovery}.
\end{align*}
The current evidence suggests that this order is essential.  The homotopy pass
finds the correct basins; resultants, multivariate sweeps, and projective charts
repair the missing sheets; and every candidate is ultimately judged by the same
Pandrosion polish on the original complex system.

\begin{thebibliography}{99}

\bibitem{smale17}
S.~Smale,
\emph{Mathematical problems for the next century},
The Mathematical Intelligencer \textbf{20} (1998), 7--15.

\bibitem{beltranpardo}
C.~Beltran and L.~M.~Pardo,
\emph{Smale's 17th problem: average polynomial time to compute affine and projective solutions},
Journal of the American Mathematical Society \textbf{22} (2009), 363--385.

\bibitem{lairez}
P.~Lairez,
\emph{A deterministic algorithm to compute approximate roots of polynomial systems in polynomial average time},
Foundations of Computational Mathematics \textbf{17} (2017), 1265--1292.

\bibitem{schmidt}
J.~W.~Schmidt,
\emph{Eine Ubertragung der Regula Falsi auf Gleichungen in Banachraumen},
Zeitschrift fur Angewandte Mathematik und Mechanik \textbf{43} (1963), 1--8.

\bibitem{ortega}
J.~M.~Ortega and W.~C.~Rheinboldt,
\emph{Iterative Solution of Nonlinear Equations in Several Variables},
Academic Press, 1970.

\bibitem{shubsmale}
M.~Shub and S.~Smale,
\emph{Complexity of Bezout's theorem I: geometric aspects},
Journal of the American Mathematical Society \textbf{6} (1993), 459--501.

\bibitem{hss}
J.~H.~Hubbard, D.~Schleicher, and S.~Sutherland,
\emph{How to find all roots of complex polynomials by Newton's method},
Inventiones Mathematicae \textbf{146} (2001), 1--33.

\bibitem{steffensen}
J.~F.~Steffensen,
\emph{Remarks on iteration},
Skandinavisk Aktuarietidskrift \textbf{16} (1933), 64--72.

\end{thebibliography}

\end{document}
